\chapter{Examples, counterexamples and failed shortcuts}\label{ch:examples}
\section{Small exact collisions}
The following identities are mandatory regression tests:
\[
 h(12)=h(14)=336,\qquad
 h(315)=h(351)=196560,\qquad
 h(160)=h(189)=60480.
\]
The first disproves cubefree injectivity and injectivity on the union of all
two-prime supports; it does not contradict fixed-pair rigidity. The second
uses odd inputs but an even target. The third has coprime inputs, so distinct
preimages need not share a prime factor.

\section{The 29-digit odd cubefree collision}\label{examples:M}
Define
\begin{align*}
 a={}&11^2 13^2\cdot17\cdot19\cdot31\cdot61\cdot97\cdot127^2
       \cdot271\cdot307^2\cdot331\cdot367,\\
 b={}&13\cdot17^2 19^2\cdot23\cdot31^2\cdot43\cdot61^2
       \cdot83\cdot127\cdot307\cdot733\cdot5419.
\end{align*}
Then
\[
 a=60629697601617236747379985403,\qquad
 b=61655391632564660609643619057,
\]
\[
 M=h(a)=h(b)=5276516179938729922490847708056660616574496169354744299520.
\]
Its factorization is
\begin{align*}
 M={}&2^{20}3^5\cdot5\cdot7^3\cdot11^2 13^2 17^2 19^2
       \cdot23\cdot31^2\cdot43\cdot61^2\\
    &{}\cdot83\cdot97\cdot127^2\cdot271\cdot307^2\cdot331\cdot367
       \cdot733\cdot5419.
\end{align*}
Every input prime is at least $11$. No prime has the same positive exponent
in both inputs. Their combined support has 17 bases. Each input has 20,736
divisors; complete divisor summation gives
\[
 \sigma(a)=87028574917351791526375587840,\qquad
 \sigma(b)=85580774693381738895491727360.
\]
The exact inverse programs give $f(M)=f_2(M)=2$.
Because $\gcd(ab,42)=1$, the four numbers $12a,14a,12b,14b$ are distinct
cubefree preimages of $336M$. Exact enumeration gives
\[
 f_2(336M)=4,\qquad f(336M)=6.
\]
The two additional unrestricted preimages are not cubefree. Their full values
are preserved in the branch-solver results and consolidated machine-readable
catalogue in the handoff archive.

This example invalidates odd-cubefree injectivity, cubefree multiplicity bounds
of two or three, and the conjecture that every cubefree collision is a
common-unitary scaling of $(12,14)$. It does not disprove a larger absolute
cubefree multiplicity bound. The earlier scan through input five million was
consistent with those false guesses only because this witness lay far beyond
that range. Its support was found by valuation constraints rather than by a
consecutive-input scan.

The exact denominator audit gives
\[
 R=2^{19}3^2\cdot5\cdot7^2\cdot11\cdot31,
 \qquad T=2^{20}3^5\cdot5\cdot7^3,\qquad T/R=378/341.
\]
Sixteen small-prime labels change, even though the complete exact solver needs
one binary branch at prime $11$. Changed coordinates and independent choices
are not the same quantity.

\section{The singleton fractional-obstruction target}
\[
 N_0=504654433920=2^7 3^3\cdot5\cdot7^2\cdot13\cdot19^2\cdot127,
\]
\[
 h^{-1}(N_0)=\{492765\},\qquad
 492765=3\cdot5\cdot7\cdot13\cdot19^2.
\]
Its cap-two system after row-support propagation has 17 columns, rank 14,
and a strictly positive fractional feasible point. The full parameterization
and the three integer deductions forcing its unique solution appear in
Chapter~\ref{ch:integer}. A complete unrestricted check evaluates all 2,304
target divisors. This is the test to apply to any claim that zero-budget linear
certificates always remove fractional freedom.

\section{Other complete fibers}
The exact six-, seven- and nine-preimage targets are
\begin{align*}
 N_6&=7089671638182002688000
     =2^{31}3^2 5^3\cdot7\cdot13\cdot31\cdot127\cdot8191,\\
 N_7&=106345074572730040320
     =2^{28}3^3\cdot5\cdot7\cdot13\cdot31\cdot127\cdot8191,\\
 N_9&=1826980530660612389572800675840\\
    &=2^{45}3^3\cdot5\cdot7\cdot13\cdot31\cdot127\cdot8191\cdot131071.
\end{align*}
All inputs are listed in the machine-readable catalogue. The full seven-member
fiber and its exchange dictionary are also in Chapter~\ref{ch:exchanges}.
Their unrestricted verifiers examine 12,288, 7,424 and 23,552 divisors,
respectively. These particular-fiber checks are not reproductions of a global
minimality census in the external paper.

The target $N_1=h(37^5 43^3)$ from Chapter~\ref{ch:bh} equals
\[
 31984991189820195724150800.
\]
Its full fiber is the singleton $\{5513329989199\}$, verified by testing all
33,600 target divisors. Five polynomial-pair target constructions, using the
first $1,2,3,5,8$ pairs, have unique cap-five preimages, despite their unpruned
positive cycles. Their targets have 26, 63, 105, 204 and 391 decimal digits.
The larger four checks are cap-five only, not unrestricted censuses.

\section{Local divisibility obstructions}
The identities
\[
 \sigma(29^2)=13\cdot67,
 \qquad \sigma(269^2)=13\cdot37\cdot151
\]
serve different purposes. In the first, $67\equiv1\pmod3$ but is $1$ modulo
neither $5$ nor $7$, while $\sigma(67^4)$ and $\sigma(67^6)$ are coprime to
$2\cdot3\cdot5\cdot7$. A third exponent-index type can therefore supply a
prime involved in an incomparable exchange of the other two types. This
blocks the direct extension of the two-index residue proof.
In the second, every factor exceeds $7$ and is smaller than $269$.
Rough divisor sums need not introduce a larger prime.

Likewise, $7$ is primitive at index three for both bases $2$ and $11$.
A primitive prime is not automatically new globally, and need not exceed its
base. The valid local distinctness in Chapter~\ref{ch:primitive} must not be
strengthened silently.

\section{Long graph cycles are not arithmetic collisions}
The following eight-prime cycle is entirely quadratic:
\[
 193\to1783\to1279\to2383\to12541\to576151\to151009\to6067\to193.
\]
For consecutive vertices $(p,q)$, the integer quotients
$(p^2+p+1)/q$ are, in order,
\[
 21,\ 2487,\ 687,\ 453,\ 273,\ 2198217,\ 3758673,\ 190749.
\]
Its eight negative edges form a positive signed cycle. It does not by itself
supply two solutions of the full valuation equations.
The nine-vertex positive cycle in Chapter~\ref{ch:universal} has all vertices
larger than the square root of its largest vertex. It refutes that particular
small-vertex deletion rule, not an asymptotic bound for actual fibers.

\section{Failure map: statements not to reintroduce}
\begin{longtable}{@{}p{0.38\textwidth}p{0.55\textwidth}@{}}
\toprule Proposed shortcut & Why it fails or remains unproved\\\midrule\endhead
Count powerful $d$ with $h(d)\mid N$ & The odd represented family in Chapter~\ref{ch:foundations} has a positive-constant relaxed count. Exact squarefree completion matters.\\
Optimize $\prod\alpha_p$ to constant zero & Its own maximal-order constant is $\log3/3$. This is not a lower bound for $f$.\\
More prime-index entropy weights suffice & Uniform exponents on $\{0,1,2\}$ saturate the local relaxation at $C_0$. Cross-prime constraints are omitted.\\
Comparable-index restriction alone suffices & Abstract exponentially large exponent families satisfy its incompatibility and aggregate size constraints.\\
Drop compatibility in the small-multiplier theorem & The multiplier becomes $T/R$, with uncontrolled denominator; the 29-digit collision changes 16 labels.\\
Odd inputs behave like odd targets & False; the large collision uses odd cubefree inputs and an even target.\\
Real nullity measures genuine choices & A singleton fiber has a 3-dimensional positive real relaxation. Integer rounding is essential.\\
Positive cycles automatically create collisions & The polynomial prime pairs have such cycles but fixed-pair injectivity holds for all exponents.\\
Universal feedback cost is surely subpolynomial & Conditional Bateman--Horn lower bounds obstruct this for caps at least five. This is not an unconditional disproof.\\
Finite graph transversals imply slow asymptotic growth & Cutoffs through $10^6$ are finite data, not an asymptotic estimate.\\
Minimum Hamming distance three suffices & Abstract independent $000/111$ triples form exponentially many words. Actual arithmetic constraints are still required.\\
Exchange decomposition already bounds multiplicity & It does not bound the number, packing, or arithmetic costs of exchanges. Extracting exchanges from a previously enumerated fiber is circular as a uniform bound.\\
Fixed-abundancy counts transfer directly & $\sigma(k)/k=N/k^2$ varies throughout a fiber.\\
A reported $x^{o(1)}$ theorem has the needed scale & Its exponent may still be a positive multiple of $\log x/\log\log x$, or even larger. Inspect the quantitative proof.\\
\bottomrule
\end{longtable}

\section{A second large witness, specified without decimal ambiguity}
The following table defines $a=\prod p^{e_p}$ and $b=\prod p^{f_p}$.
They are odd cubefree integers with $h(a)=h(b)$, each input has 106 decimal
digits, and the common target has 211. The 47-base support lies between 17 and
9439. Exact geometric products and independently certified valuation products
were checked; the full divisor sets were not enumerated. This is a witness,
not a complete-fiber or minimality claim.
\begin{center}\small
\begin{tabular}{rrr@{\qquad}rrr@{\qquad}rrr}
\toprule
$p$&$e_p$&$f_p$&$p$&$e_p$&$f_p$&$p$&$e_p$&$f_p$\\\midrule
17&1&0&103&1&2&431&0&1\\
19&1&2&109&0&2&433&1&0\\
23&0&1&127&2&1&499&2&1\\
31&1&2&137&1&2&547&1&2\\
37&1&2&139&1&2&571&2&1\\
41&0&2&151&1&2&613&1&0\\
47&2&1&157&0&1&733&2&1\\
53&1&2&163&1&0&1009&1&0\\
59&1&0&181&0&1&1093&1&0\\
61&1&2&271&0&1&1723&1&0\\
67&1&0&307&1&2&3169&0&1\\
73&1&0&313&1&0&3463&0&1\\
79&2&1&317&1&0&3571&2&1\\
83&1&2&331&1&0&5419&1&2\\
97&2&1&367&2&0&9439&0&1\\
101&0&1&409&1&0&&&\\
\bottomrule
\end{tabular}
\end{center}
The exact decimal inputs, target, and signed-difference columns are in \texttt{general\_17\_10000\_1.json} in the v5 packet. They can also be reconstructed uniquely from this table.
